\documentclass[10pt]{article}
\usepackage[T1]{fontenc}
\usepackage[utf8]{inputenc}
\usepackage[french]{babel}
\usepackage{amsmath,amssymb,mathrsfs,mathpazo}
\usepackage{geometry}
\geometry{a4paper,margin=2cm}
\usepackage{graphicx,array,tabularx,tabularray,colortbl,multirow,multicol,booktabs,enumitem,pifont,ulem,fancybox,diagbox,cancel}
\usepackage{pstricks,pstricks-add,pst-plot,pst-tree,pst-node,pst-text,pst-func,pst-math,pst-bspline,pst-3dplot,pst-eucl,pst-solides3d}
\let\pscalculate\undefined
\usepackage{pst-bezier}
\usepackage{tikz,pgfplots,tkz-tab}
\usetikzlibrary{arrows,arrows.meta,patterns,calc,positioning,trees,decorations.markings,shapes.geometric,shapes.misc,shapes.arrows,intersections,decorations.pathmorphing,decorations.pathreplacing,backgrounds,shadows,babel}
\pgfplotsset{compat=1.18}
\usepgfplotslibrary{fillbetween}
\usepackage[np]{numprint}
\usepackage{listings,makecell,pas-tableur,twoopt}
\let\barre\undefined
\usepackage[pstricks]{tablvar}
\definecolor{darkgreen}{RGB}{0,100,0}
\providecommand{\red}{\color{red}}
\providecommand{\blue}{\color{blue}}
\providecommand{\green}{\color{green}}
\providecommand{\black}{\color{black}}
\providecommand{\white}{\color{white}}
\providecommand{\gray}{\color{gray}}
\providecommand{\cyan}{\color{cyan}}
\providecommand{\magenta}{\color{magenta}}
\providecommand{\yellow}{\color{yellow}}
\providecommand{\violet}{\color{violet}}
\providecommand{\orange}{\color{orange}}
\providecommand{\euro}{\texteuro}
\providecommand{\EUR}{\texteuro}
\providecommand{\ieme}{\textsuperscript{e}}
\providecommand{\degres}{\ensuremath{^\circ}}
\providecommand{\degree}{\ensuremath{^\circ}}
\providecommand{\vv}[1]{\overrightarrow{#1}}
\providecommand{\Coord}[1]{#1}
\providecommand{\decosix}{}
\providecommand{\starredbullet}{\ensuremath{\bullet}}
\pagestyle{empty}
\setlength{\parindent}{0pt}
\usepackage{adjustbox}\DeclareUnicodeCharacter{2236}{:}\DeclareUnicodeCharacter{2190}{\ensuremath{\leftarrow}}
\providecommand{\R}{\mathbb{R}}
\providecommand{\N}{\mathbb{N}}
\providecommand{\Z}{\mathbb{Z}}
\providecommand{\Q}{\mathbb{Q}}
\providecommand{\D}{\mathbb{D}}
\providecommand{\C}{\mathbb{C}}
\providecommand{\e}{{\mathrm{e}}}
\providecommand{\vect}[1]{\overrightarrow{#1}}
\providecommand{\vectt}[1]{\overrightarrow{\mathrm{#1}}}
\providecommand{\vvt}[1]{\overrightarrow{\mathrm{#1}}}
\providecommand{\np}[2][]{{#2~#1}}
\providecommand{\vv}[1]{{\overrightarrow{#1}}}
\renewcommand{\d}{{\mathrm{d}}}
\providecommand{\degres}{{^\circ}}
\providecommand{\degre}{{^\circ}}
\providecommand{\blue}{\color{blue}}
\providecommand{\red}{\color{red}}
\providecommand{\enskip}{\quad}
\providecommand{\llbracket}{\lbrack\!\lbrack}
\providecommand{\rrbracket}{\rbrack\!\rbrack}
\providecommand{\euro}{{\text{€}}}
\providecommand{\no}{{\text{n°}}}
\providecommand{\interval}[2]{{[#1;#2]}}
\providecommand{\intervFO}[2]{{[#1;#2[}}
\providecommand{\dotfill}{\ldots}

\providecommand{\ds}{\displaystyle}
\providecommand{\textding}[1]{#1}
\providecommand{\Oij}{\(\mathrm{O};\vec{\imath},\vec{\jmath}\)}
\providecommand{\Oijk}{\(\mathrm{O};\vec{\imath},\vec{\jmath},\vec{k}\)}
\providecommand{\Ouv}{\(\mathrm{O};\vec{u},\vec{v}\)}
\providecommand{\intervFF}[2]{\ensuremath{[#1;#2]}}
\renewcommand{\euro}{\ifmmode\text{\texteuro}\else\texteuro\fi}
\DeclareUnicodeCharacter{00B0}{\ifmmode^{\circ}\else\textdegree\fi}
\DeclareUnicodeCharacter{20AC}{\ifmmode\text{\texteuro}\else\texteuro\fi}
% Énoncés : Éducation nationale. Transcription et corrigés : APMEP, Denis Vergès et contributeurs. Découpage et classement : CoopMaths, Sébastien Lozano, Éric Elter et contributeurs. Adaptation de la présentation : Maths Explorer.
% CC BY-SA : https://coopmaths.fr/www/about/
% https://coopmaths.fr/alea/static/e3c/2020/tex/e3c_2020_00_sujet25_2.tex
\begin{document}
\noindent\textit{APMEP / CoopMaths. CC BY-SA. Présentation adaptée par Maths Explorer.}\par\medskip
\medskip

Le principe d'un Escape Game est le suivant : une équipe de participants est enfermée à l'intérieur d'une salle à thème et doit réussir à en sortir en moins d'une heure (on parle alors de partie réussie). Au-delà d'une heure, les participants sont libérés et la partie est perdue.

Un exploitant d'Escape Game propose à ses participants de faire deux parties à la suite : la première partie se déroule dans la salle à thème \og Espion \fg, la seconde partie dans la salle à thème \og Musée \fg{}. Il dispose des données suivantes :
\begin{itemize}
\item lorsqu'une équipe joue dans la salle à thème \og Espion \fg{}, la probabilité qu'elle réussisse sa partie \og Espion \fg{} est égale à 0,5 ;
\item lorsqu'une équipe a réussi la partie \og Espion\fg, la probabilité qu'elle réussisse sa partie \og Musée \fg{} est égale à 0,6 ;
\item lorsqu'une équipe n'a pas réussi la partie \og Espion \fg, la probabilité qu'elle réussisse sa partie \og Musée \fg{} est égale à 0,45.
\end{itemize}
 
Une équipe est choisie au hasard. On note les événements suivants :
\begin{itemize}
\item $E$ : \og l'équipe réussit la partie \og Espion \fg{} ;
\item $M$ : \og l'équipe réussit la partie \og Musée \fg.
\end{itemize} 

\medskip

\begin{enumerate}
\item Sur la copie, recopier et compléter l'arbre de probabilités suivant :

\begin{center}
\psset{nodesepA=0pt,nodesepB=3pt,treesep=0.75,labelsep=0.1pt,levelsep=2.5cm}
\pstree[treemode=R]{\TR{}}
{\pstree{\TR{$E$~~}\taput{$\dots$}}
	{
	\TR{$M$}\taput{$\dots$}
	\TR{$\overline{M}$}\tbput{$\dots$}
	}
\pstree{\TR{$\overline{E}$~~}\tbput{$\dots$}}
	{\TR{$M$}\taput{$\dots$}
	\TR{$\overline{M}$}\tbput{$\dots$}
	}
}
\end{center}

\item Déterminer la probabilité que l'équipe réussisse les deux parties.
\item Montrer que la probabilité que l'équipe réussisse la partie \og Musée \fg{} est égale à 0,525.
\item Quelle est la probabilité qu'une équipe échoue à la partie \og Espion \fg{} sachant qu'elle a réussi la partie \og Musée \fg{} ? On donnera la réponse arrondie à $10^{-2}$.
\item Pour chacune des deux parties qui sont gagnées, une équipe reçoit 2 \euro{} de réduction pour une prochaine visite. Elle peut donc recevoir 0, 2 ou 4 \euro{} de réduction.

Si un très grand nombre d'équipes jouent les deux parties, quel est le montant moyen de la réduction obtenue à la fin des deux parties ? Expliquer la démarche.
\end{enumerate}

\vspace{0,5cm}

\end{document}
